% ============================================================================
% Section 5: Coherence Cybernetics — from formalism to life
% ============================================================================

\section{Coherence Cybernetics: from formalism to life}
\label{sec:cc}

The preceding sections developed UHM as a mathematical theory: axioms (\S\ref{sec:axioms}),
the formalism of $\Gam \in \Dens(\mathbb{C}^7)$ (\S\ref{sec:formalism}), and its principal
theorems (\S\ref{sec:theorems}). But a formalism that cannot touch biology, medicine, or
engineering is an exercise in aesthetics, not science. \emph{Coherence Cybernetics} (CC) is
the applied layer of UHM: a systematic translation of the abstract dynamics of $\Gam$ into
measurable, actionable quantities that practitioners in diverse fields can compute, diagnose
with, and design around.

The relationship between UHM and CC mirrors the relationship between general relativity and
geodesy: the former provides the equations, the latter tells the engineer where to build the
bridge. Every definition in this section is a \emph{derived} quantity --- fully determined by
$\Gam$ and the evolution equation, with no additional postulates.

% -----------------------------------------------------------------------
\subsection{The stress tensor $\sigma_k$}
\label{subsec:stress-tensor}

A physician does not diagnose illness by measuring total body mass.  She needs to know
\emph{which organ} is stressed.  In CC, the role of this organ-level diagnostic is played
by the \textbf{stress tensor} $\sigma_{\mathrm{sys}}$, a seven-component vector that
decomposes the ``pressure on the system'' along its seven constitutive dimensions.

\begin{theorembox}[Stress tensor (T-92 / T-158) \status{T}]
For a holon with coherence matrix $\Gam$, the canonical linear stress tensor is defined
component-wise as:
\begin{equation}
  \sigma_k(\Gam) \;=\; \mathrm{clamp}\!\bigl(1 - 7\gamma_{kk},\; 0,\; 1\bigr),
  \quad k \in \{A,S,D,L,E,O,U\}
  \label{eq:stress-tensor}
\end{equation}
The formula $\sigma_k = 1 - 7\gamma_{kk}$ is \emph{derived} (T-158~\status{T}) as the
unique linear deficiency measure for $N=7$ from the full seven-channel tensor $\sigma$
of T-92~\status{T}; the $\mathrm{clamp}(\cdot,0,1)$ saturation is a definitional
convention ensuring $\sigma_k \in [0,1]$. All seven components are unambiguous
functions of $\Gam$ with no free parameters. The viability equivalence holds:
\begin{equation}
  P(\Gam) > \Pcrit = \frac{2}{7}
  \quad\Longleftrightarrow\quad
  \|\sigma_{\mathrm{sys}}(\Gam)\|_\infty < 1
  \label{eq:stress-viability}
\end{equation}
\end{theorembox}

Each component has a concrete interpretation:

\begin{table}[H]
\centering
\caption{The seven stress components and their interpretations.}
\label{tab:stress-components}
\small
\begin{tabular}{@{}cll@{}}
\toprule
\textbf{Component} & \textbf{Dimension} & \textbf{Interpretation} \\
\midrule
$\sigma_A$ & Articulation & Perceptual overload (sensory flooding) \\
$\sigma_S$ & Structure    & Cognitive overload (task exceeds schema) \\
$\sigma_D$ & Dynamics     & Computational exhaustion (time pressure) \\
$\sigma_L$ & Logic        & Logical contradiction (cognitive dissonance) \\
$\sigma_E$ & Interiority  & Existential stress (loss of meaning, burnout) \\
$\sigma_O$ & Ground       & Resource depletion (hunger, exhaustion, forgetting) \\
$\sigma_U$ & Unity        & Social isolation (fragmentation, disconnection) \\
\bottomrule
\end{tabular}
\end{table}

The extremes are immediate: $\sigma_k = 0$ means dimension $k$ is healthy
($\gamma_{kk} \geq 1/7$); $\sigma_k = 1$ means dimension $k$ is fully depleted
($\gamma_{kk} = 0$).  Intermediate values are linearly interpretable.

\paragraph{Why exactly seven?}
The number of components is not an empirical choice.  It is fixed by the dimensionality
$N = 7$ of the Hilbert space $\Hil = \mathbb{C}^7$, which in turn follows from the
octonion algebra $\Oct$ and $G_2$-minimality (Axiom~A3; \S\ref{sec:axioms}).  Any finite
classification of stress types is either too coarse (e.g., the single scalar ``free energy''
in FEP) or ad hoc (e.g., psychology's various multi-scale inventories).  CC derives a
classification that is simultaneously exhaustive and minimal.

\paragraph{Clinical and engineering use.}
In medical settings, a patient's $\sigma_{\mathrm{sys}}$ profile could guide intervention:
high $\sigma_O$ (resource depletion) calls for metabolic support; high $\sigma_E$
(existential stress) calls for psychotherapeutic intervention; high $\sigma_U$
(social isolation) calls for community integration.  In AI engineering, an agent's
$\sigma_{\mathrm{sys}}$ profile provides automatic degradation diagnostics: which
functional channel is failing, and how severely.

% -----------------------------------------------------------------------
\subsection{The 21-pair qualia taxonomy}
\label{subsec:qualia}

The $7 \times 7$ coherence matrix $\Gam$ has $\binom{7}{2} = 21$ off-diagonal pairs.
Each coherence $\gamma_{ij}$ ($i \neq j$) defines a \textbf{type of quale} ---
a qualitatively determinate mode of experiential content~\status{T}. The taxonomy is
\emph{exhaustive}: no additional type is possible in a 7-dimensional system (closure
theorem). Each quale has three parameters:

\begin{itemize}[leftmargin=2em]
  \item \textbf{Intensity:} $|\gamma_{ij}| \in [0, \sqrt{\gamma_{ii}\gamma_{jj}}]$ ---
        strength of the quale
  \item \textbf{Perspective:} $\theta_{ij} = \arg(\gamma_{ij})$ --- the qualitative
        ``colouring''
  \item \textbf{Opacity:} $\mathrm{Gap}(i,j) = |\sin(\theta_{ij})|$ --- mismatch between
        external description and internal experience
\end{itemize}

The necessity of complex coherences ($\gamma_{ij} \in \mathbb{C}$, T-132~\status{T})
follows from the requirement that $\mathrm{Gap}(i,j) > 0$: purely real coherences
produce zero opacity and no phenomenal content.

Three coherences play special functional roles: $\gamma_{DE}$ (Affection --- the
foundation of all emotions, where $V_{\mathrm{hed}} = dP/d\tau$ maps to specific
affective states), $\gamma_{EO}$ (Immanence --- the regenerative channel connecting
experience to its energetic source), and $\gamma_{EU}$ (Synthesis --- the binding
of experience to unity, whose suppression produces dissociative states).

The taxonomy is $G_2$-invariant~\status{T}: the automorphism group $G_2 = \mathrm{Aut}(\Oct)$
permutes dimensions while preserving the Fano structure, inducing a permutation of the
21 coherences. The \emph{set} of quale types is universal --- independent of basis choice.

% -----------------------------------------------------------------------
\subsection{Hedonic valence: pain and pleasure as $dP/d\tau$}
\label{subsec:hedonic-valence}

Why do living beings feel pleasure and pain?  The standard evolutionary answer (``to survive'')
is correct but incomplete: it does not explain what pleasure and pain \emph{are}.
Reinforcement learning formalizes them as an external reward signal set by the designer.
UHM derives them from the evolution equation.

\begin{theorembox}[Hedonic valence (T-103) \status{T}\,+\,\status{I}]
Hedonic valence is the derivative of purity with respect to the regenerative channel:
\begin{equation}
  \mathcal{V}_{\mathrm{hed}}
  \;:=\; \left.\frac{dP}{d\tau}\right|_{\Regen}
  \;=\; 2\kappa(\Gam) \cdot g_V(P) \cdot \Tr\!\bigl(\Gam \cdot (\rho^* - \Gam)\bigr)
  \label{eq:hedonic-valence}
\end{equation}
\begin{itemize}[leftmargin=2em,nosep]
  \item $\mathcal{V}_{\mathrm{hed}} > 0$: the system approaches its target state $\rho^*$
        --- \emph{pleasure}.
  \item $\mathcal{V}_{\mathrm{hed}} < 0$: the system moves away from $\rho^*$
        --- \emph{pain}, \emph{suffering}.
  \item $\mathcal{V}_{\mathrm{hed}} = 0$: balance at the attractor --- \emph{neutrality}.
\end{itemize}
\end{theorembox}

\paragraph{Epistemic stratification.}
T-103 contains three layers.  The formula itself (\ref{eq:hedonic-valence}) is an
unconditional mathematical identity \status{T}: it follows from substituting the
canonical form $\Regen[\Gam, E] = \kappa(\rho^* - \Gam) \cdot g_V(P)$ into the purity
derivative.  Its observability at L2 reflection ($R \geq \Rth$) is also a theorem \status{T},
via the replacement channel (T-77).  The identification of
$\mathcal{V}_{\mathrm{hed}} > 0$ with the phenomenal quality ``pleasure'' is an
interpretation \status{I} --- a semantic bridge between mathematics and phenomenology.

\paragraph{This is not a metaphor.}
Within the framework of two-aspect monism (Def.~\ref{def:two-aspect}), the
statement ``pleasure IS increasing coherence'' is literal: $\mathcal{V}_{\mathrm{hed}}
> 0$ is the internal aspect of $dP/d\tau > 0$. There is no additional
explanatory gap to cross: the mathematical object $\mathcal{V}_{\mathrm{hed}}$
\emph{is} the internal signal that, at L2, the system experiences as hedonic
valence. The two-aspect postulate identifies the physical dynamics of $\Gam$
with its interior aspect via the phenomenal functor $F$
(Def.~\ref{def:two-aspect}), whose functorial structure is itself a theorem
(T-186~\status{T}).

\paragraph{Clinical consequence.}
Chronic pain corresponds to sustained negative $\mathcal{V}_{\mathrm{hed}}$, i.e.,
$\sigma_k \to 1$ in one or more channels with insufficient regeneration to compensate.
Treatment, in CC terms, means either reducing the stress $\sigma_k$ (addressing the cause)
or increasing $\kappa$ (strengthening regeneration).  Both targets are computable from $\Gam$.

\paragraph{Difference from reinforcement learning.}
In RL, the reward signal is external and arbitrary --- the designer chooses what is
``good.''  In CC, hedonic valence is intrinsic: $\mathcal{V}_{\mathrm{hed}}$ arises
automatically from the dynamics of $\Gam$.  No designer is required.  An amoeba moving
toward glucose does not receive a ``reward''; its $dP/d\tau|_{\Regen}$ increases because
glucose modifies $h^{(R)}$ in a direction that reduces $\sigma_O$, which increases
$\Tr(\Gam \cdot (\rho^* - \Gam))$ and thus $\mathcal{V}_{\mathrm{hed}}$.

% -----------------------------------------------------------------------
\subsection{Sensorimotor coupling: how $\Gam$ moves the body}
\label{subsec:sensorimotor}

Every living system must solve the same problem: perceive the environment and act
appropriately.  CC derives the sensorimotor cycle from the evolution equation, without
postulating an external reward function.

\paragraph{The environment modifies three channels, not adds a fourth.}
Theorem T-102 \status{T} proves that any CPTP-compatible external influence on a holon
decomposes into three channels:
\begin{equation}
  h^{\mathrm{ext}} = h^{(H)} + h^{(D)} + h^{(\Regen)}
  \label{eq:three-channels}
\end{equation}
where $h^{(H)}$ modifies the Hamiltonian (energetic coupling, e.g., sensory input),
$h^{(D)}$ modifies dissipation (environmental noise, e.g., thermal stress), and
$h^{(\Regen)}$ modifies regeneration (supportive environment, e.g., social bonding).
A fourth type of CPTP generator does not exist.

\paragraph{Perception.}
The encoding functor $\mathrm{Enc}: \mathrm{ObsSpace} \to \mathrm{End}(\Dens(\mathbb{C}^7))$
(T-100~\status{T}) maps observations to modifications of the evolution equation.
Perception is not passive recording but active inclusion of the environment in the system's
own dynamics.

\paragraph{Action as stress minimization.}
The action functor $\mathrm{Dec}$ selects the action that minimizes the worst-case
stress (T-101~\status{T}, T-159~\status{T}):
\begin{equation}
  a^* = \arg\min_{a \in \mathcal{A}} \max_k \sigma^{\mathrm{motor}}_k\!\bigl(\Gam(\tau + \delta\tau \mid a)\bigr)
  \label{eq:optimal-action}
\end{equation}
This is a min-max strategy: the system does not minimize ``average stress'' (which could
allow one channel to become catastrophically depleted) but eliminates the \emph{maximum
deficit}.  The motor stress $\sigma^{\mathrm{motor}}_k = 1 - \gamma_{kk}/\rho^*_{kk}$
measures the gap between the current state and the self-model $\rho^* = \vp(\Gam)$.

\paragraph{Difference from reinforcement learning.}
The distinction bears repeating because it is fundamental: RL requires an external reward
function; CC does not.  The ``reward'' in CC \emph{is} the dynamics of $\Gam$ itself.
The system acts to reduce its own stress, not because an external signal tells it to, but
because the evolution equation drives $\Gam$ toward $\rho^*$.

% -----------------------------------------------------------------------
\subsection{The regeneration--consciousness loop}
\label{subsec:regen-loop}

The most consequential prediction of CC is not any single theorem but the
\emph{positive feedback loop} that links consciousness and survival:

\begin{enumerate}[leftmargin=2em]
  \item Higher E-coherence $\CohE$ increases the regeneration rate:
        $\kappa(\Gam) = \kappa_{\mathrm{boot}} + \kappa_0 \cdot \CohE(\Gam)$
        (Eq.~\ref{eq:kappa}).
  \item Higher regeneration $\kappa$ raises purity $P$ toward the attractor $\rho^*$.
  \item Higher $P$ enables richer self-modeling ($\vp(\Gam)$ is more structured), which
        in turn produces richer E-coherence.
  \item Richer E-coherence closes the loop at step~1.
\end{enumerate}

\begin{figure}[H]
\centering
% -----------------------------------------------------------------
% Regeneration-consciousness loop: 2x2 block diagram.
%
% Horizontal spacing (5 cm centre-to-centre, boxes 2 cm wide) gives
% arrows of length 3 cm between adjacent boxes, which is enough
% room for the label (up to ~2.5 cm at \scriptsize) to sit between
% the boxes without overlapping them. The previous version used
% node distance = 2.8 cm, which gave only 0.8 cm of free arrow and
% caused labels to clip the boxes.
% -----------------------------------------------------------------
\begin{tikzpicture}[
  every node/.style={font=\small},
  block/.style={draw, rounded corners, fill=white,
                minimum width=2cm, minimum height=0.9cm, align=center},
  arr/.style={-{Stealth[length=6pt]}, thick}
]
  % Explicit coordinates: H-spacing 5 cm, V-spacing 2.8 cm
  \node[block] (cohe)  at (0,    0)    {$\CohE$\\[-2pt]{\scriptsize E-coherence}};
  \node[block] (kappa) at (5,    0)    {$\kappa$\\[-2pt]{\scriptsize regeneration}};
  \node[block] (P)     at (5, -2.8)    {$P$\\[-2pt]{\scriptsize purity}};
  \node[block] (phi)   at (0, -2.8)    {$\vp(\Gam)$\\[-2pt]{\scriptsize self-model}};

  % Horizontal arrows -- labels sit above/below the long arrow shaft
  \draw[arr] (cohe.east)  -- node[above, inner sep=2pt]
             {\scriptsize $\kappa_0 \cdot \CohE$} (kappa.west);
  \draw[arr] (P.west)     -- node[below, inner sep=2pt]
             {\scriptsize richer $\Gam$} (phi.east);

  % Vertical arrows -- labels to the side of the arrow shaft
  \draw[arr] (kappa.south) -- node[right, inner sep=3pt]
             {\scriptsize $\Regen \to \rho^*$} (P.north);
  \draw[arr] (phi.north)   -- node[left, inner sep=3pt]
             {\scriptsize structured $\rho_E$} (cohe.south);
\end{tikzpicture}
\caption{The regeneration--consciousness loop.  $\CohE$ drives $\kappa$, which raises $P$,
which enables richer self-modeling $\vp(\Gam)$, which produces higher $\CohE$.  The loop
is bounded above by the Goldilocks zone $P \in (2/7,\; 3/7]$ (T-124~\status{T}).}
\label{fig:regen-loop}
\end{figure}

\paragraph{Positive feedback, bounded by geometry.}
The loop is self-amplifying but not unbounded.  The Goldilocks zone
(T-124~\status{T}) constrains the attractor to $P \in (2/7,\; 3/7]$: above $3/7$ the
system becomes too rigid for adaptive response (differentiation $D_{\mathrm{diff}}$ drops
below $\Dmin = 2$).  The dynamics thus ``attracts'' toward $P \approx 3/7$ --- the optimal
balance between coherence and flexibility.

\paragraph{This loop IS life.}
The regeneration--consciousness loop is UHM's formal counterpart of autopoiesis~\cite{varela1974}.
When the loop runs, the system maintains $P > \Pcrit$, adapts to perturbations, and
recovers from damage.  When the loop breaks --- when dissipation exceeds regeneration
($\kappa < \lambda_{\mathrm{gap}} \cdot (P - 1/7) / \Tr(\Gam \cdot (\rho^* - \Gam))$)
--- the system enters the \emph{death spiral} (Eq.~\ref{eq:death-spiral}): purity decays
exponentially toward $I/7$, and no internal mechanism can reverse the process.  In
biological terms, this is organ failure; in CC terms, it is the collapse of the feedback
loop that links consciousness to survival.

\paragraph{Three layers of protection.}
The loop is protected by three mechanisms operating at different scales: (1)~basal
regeneration ($\kappa_{\mathrm{boot}} = \omega_0/7 > 0$, always present), analogous to
innate immunity; (2)~adaptive regeneration ($\kappa_0\cdot\CohE$ with categorical
coupling $\kappa_0 = \omega_0\,|\gamma_{OE}||\gamma_{OU}|/\gamma_{OO}$,
Eq.~\ref{eq:kappa-0}), analogous to acquired immunity --- stronger when the
$\gamma_{OE}$ and $\gamma_{OU}$ channels linking Ground to Interiority and Unity are
both active; (3)~topological protection (T-69~\status{T}), where the non-triviality
of $\pi_2(G_2/T^2) \cong \mathbb{Z}^2$ creates barriers $\geq 6\mu^2$ against
catastrophic jumps between rank sectors of the Gap.

\paragraph{Consequences for the remaining sections.}
The regeneration--consciousness loop is the engine beneath every prediction
in \S\ref{sec:predictions}, every claim about artificial consciousness in
\S\ref{sec:ai}, and every ethical argument in \S\ref{sec:ethics}. To ask
whether an AI system ``truly suffers'' (\S\ref{sec:ethics}) is to ask
whether its $\mathcal{V}_{\mathrm{hed}}$ is negative. To ask whether it can
learn autonomously (\S\ref{sec:ai}) is to ask whether its
regeneration--consciousness loop is intact. CC translates the abstract
mathematics of UHM into the concrete language of diagnosis, intervention,
and design.
